\section{Рабочий проект}
\subsection{Спецификация компонентов и классов системы}

\subsubsection{Спецификация класса UserApiController}
UserApiController (Компонент управления пользователями через API)

Описание: Управляет операциями, связанными с регистрацией пользователей и получением информации о текущем авторизованном пользователе через API. Взаимодействует с сущностью User и использует формы для обработки данных.

Валидация данных:
\begin{itemize}
	\item в методе userRegistration() используется форма ApiRegistrationFormType для валидации входных данных, таких как email и пароль, перед их сохранением;
	\item данные, полученные из HTTP-запросов в формате JSON, проверяются на соответствие требованиям формы.
\end{itemize}

Безопасность:
\begin{itemize}
	\item пароли хэшируются с использованием UserPasswordHasherInterface перед сохранением в базе данных;
	\item доступ к информации об авторизованном пользователе осуществляется через токен безопасности (security.token{\_}storage).
\end{itemize}

Роли и права:
\begin{itemize}
	\item взаимодействует с моделью UserRoles для получения информации о ролях пользователей;
	\item пакет FOSRestBundle используется для обработки REST-запросов и формирования ответов.
\end{itemize}

Основные методы:
\begin{itemize}
	\item userRegistration(Request \$request, UserRolesRepository \$userRolesRepository, EntityManagerInterface \$entityManager, UserPasswordHasherInterface \$passwordHasher): регистрирует нового пользователя, валидирует данные, хэширует пароль и сохраняет пользователя в базе данных;
	\item getAuthUserInfo(UserRepository \$repository): возвращает информацию об авторизованном пользователе;
	\item getUserRoles(UserRolesRepository \$repository): возвращает список всех ролей пользователей.
\end{itemize}
Зависимости:
\begin{itemize}
	\item взаимодействует с сущностями User и UserRoles;
	\item использует ApiRegistrationFormType для обработки форм регистрации;
	\item зависит от EntityManagerInterface для работы с базой данных;
	\item использует FOSRestBundle для обработки REST-запросов и формирования ответов в формате JSON;
	\item применяет UserPasswordHasherInterface для безопасного хэширования паролей.
\end{itemize}

\begin{xltabular}{\textwidth}{|X|X|}
	\caption{Данные класса UserApiController \label{User:table}} \\
	\hline 
	\centrow Данные & 
	\centrow Описание \\
	\hline 
	\endfirsthead
	
	\caption*{Продолжение таблицы \ref{User:table}} \\
	\hline 
	\centrow Данные & 
	\centrow Описание  \\
	\hline 
	\endhead
	
	Request &  Данные из HTTP-запросов (JSON) для валидации и обработки форм. \\ \hline
	User &  Сущность, представляющая пользователя. \\ \hline
	UserRoles &  Сущность, представляющая роли пользователей. \\ \hline
	ApiRegistrationFormType &  Форма для валидации данных при регистрации пользователя. \\ \hline
	UserPasswordHasherInterface &  Интерфейс для хэширования паролей. \\ \hline
\end{xltabular}

\begin{xltabular}{\textwidth}{|X|X|}
	\caption{Методы класса  UserApiController \label{UserMet:table}} \\
	\hline 
	\centrow Метод & 
	\centrow Описание \\
	\hline 
	\endfirsthead
	
	\caption*{Продолжение таблицы \ref{UserMet:table}} \\
	\hline 
	\centrow Метод & 
	\centrow Описание  \\
	\hline 
	\endhead
	
	userRegistration(Request \$request, UserRolesRepository \$userRolesRepository, EntityManagerInterface \$entityManager, UserPasswordHasherInterface \$passwordHasher) &  Валидирует входные данные, хэширует пароль и сохраняет нового пользователя в базе данных. \\ \hline
	getAuthUserInfo(UserRepository \$repository) &  Возвращает информацию об авторизованном пользователе в формате JSON с кодом 200. \\ \hline
	getUserRoles(UserRolesRepository \$repository) &  Возвращает список всех ролей пользователей в формате JSON с кодом 200. \\ \hline
\end{xltabular}

\subsubsection{Спецификация класса VacancyApiController}
VacancyApiController (Компонент управления вакансиями через API)

Описание: Управляет операциями, связанными с вакансиями, такими как создание, получение списка вакансий, получение данных об отдельной вакансии и получение вакансий, связанных с авторизованным пользователем. Реализует REST API и взаимодействует с сущностью Vacancy.

Валидация данных:

\begin{itemize}
	\item в методе newVacancy() используется форма ApiVacancyType для валидации входных данных, полученных в формате JSON;
	\item валидация обеспечивает корректность полей, необходимых для создания вакансии.
\end{itemize}

Безопасность:

\begin{itemize}
	\item доступ к информации о вакансиях текущего пользователя осуществляется через метод getUser(), который возвращает данные авторизованного пользователя;
	\item используется FOSRestBundle для обработки REST-запросов и формирования стандартизированных ответов.
\end{itemize}

Роли и права:
\begin{itemize}
	\item доступ к методу userResume() ограничен для авторизованных пользователей, так как он возвращает вакансии, связанные с текущим пользователем.
\end{itemize}

Основные методы:
\begin{itemize}
	\item getManyVacancies(Request \$request, VacancyRepository \$vacancyRepository): возвращает отфильтрованный список вакансий с пагинацией и информацией о фильтрах;
	\item newVacancy(Request \$request, EntityManagerInterface \$entityManager): создает новую вакансию, валидирует данные и сохраняет их в базе данных;
	\item getSingleVacancy(Vacancy \$vacancy, VacancyRepository \$vacancyRepository): возвращает данные об одной вакансии по её идентификатору;
	\item userResume(VacancyRepository \$vacancyRepository): возвращает список вакансий, связанных с авторизованным пользователем.
\end{itemize}

Зависимости:
\begin{itemize}
	\item взаимодействует с сущностью Vacancy и репозиторием VacancyRepository;
	\item использует ApiVacancyType для обработки форм создания вакансий;
	\item зависит от EntityManagerInterface для работы с базой данных;
	\item использует FilterService для обработки и применения фильтров в методе getManyVacancies;
	\item применяет FOSRestBundle для обработки REST-запросов и формирования JSON-ответов.
\end{itemize}

\begin{xltabular}{\textwidth}{|X|X|}
	\caption{Данные класса VacancyApiController \label{Vacancy:table}} \\
	\hline 
	\centrow Данные & 
	\centrow Описание \\
	\hline 
	\endfirsthead
	
	\caption*{Продолжение таблицы \ref{Vacancy:table}} \\
	\hline 
	\centrow Данные & 
	\centrow Описание  \\
	\hline 
	\endhead
	
	Request &  Данные из HTTP-запросов (JSON) для валидации и обработки форм. \\ \hline
	Vacancy &  Сущность, представляющая вакансию. \\ \hline
	ApiVacancyType &  Форма для валидации данных при создании вакансии. \\ \hline
	FilterService &  Сервис для обработки и применения фильтров к списку вакансий. \\ \hline
	EntityManagerInterface &  Интерфейс для работы с базой данных через Doctrine. \\ \hline
\end{xltabular}

\begin{xltabular}{\textwidth}{|X|X|}
	\caption{Методы класса  VacancyApiController \label{VacancyM:table}} \\
	\hline 
	\centrow Метод & 
	\centrow Описание \\
	\hline 
	\endfirsthead
	
	\caption*{Продолжение таблицы \ref{VacancyM:table}} \\
	\hline 
	\centrow Метод & 
	\centrow Описание  \\
	\hline 
	\endhead
	
	getManyVacancies(Request \$request, VacancyRepository \$vacancyRepository) &  Возвращает отфильтрованный список вакансий с пагинацией и информацией о фильтрах в формате JSON с кодом 200. \\ \hline
	newVacancy(Request \$request, EntityManagerInterface \$entityManager) &  Валидирует входные данные, сохраняет новую вакансию в базе данных. \\ \hline
	getSingleVacancy(Vacancy \$vacancy, VacancyRepository \$vacancyRepository) &  Возвращает данные об одной вакансии по её идентификатору в формате JSON. \\ \hline
	userVacancy(VacancyRepository \$vacancyRepository) &  Возвращает список вакансий, связанных с авторизованным пользователем. \\ \hline
\end{xltabular}

\subsubsection{Спецификация класса VacancyParametersApiController}
VacancyParametersApiController (компонент управления параметрами вакансий через API)

Описание:Управляет операциями получения параметров вакансий через REST API. Обеспечивает получение списка всех параметров определенного типа или данных об одном параметре по его идентификатору. Взаимодействует с сущностями, связанными с параметрами вакансий.

Валидация данных:
\begin{itemize}
	\item отсутствует валидация входных данных, так как методы обрабатывают GET-запросы с параметрами, передаваемыми в URL;
	\item проверяется существование запрошенного типа параметра вакансии, в случае ошибки выбрасывается исключение NotFoundHttpException.
\end{itemize}

Связи:
\begin{itemize}
	\item параметры вакансий связаны с сущностями в пространстве имен App/Entity/Vacancy/Vacancy, определяемыми динамически через параметр vacancyParameter.
\end{itemize}

Безопасность:
\begin{itemize}
	\item использует FOSRestBundle для обработки REST-запросов и формирования стандартизированных JSON-ответов;
	\item не реализует дополнительной аутентификации или авторизации, доступ к методам предполагается открытым.
\end{itemize}

Основные методы:
\begin{itemize}
	\item getManyVacancyParameters(string \$vacancyParameter, ManagerRegistry \$doctrine): возвращает список всех параметров вакансии для указанного типа;
	\item getSingleVacancyParameters(string \$vacancyParameter, int \$id, ManagerRegistry \$doctrine): возвращает данные об одном параметре вакансии по его идентификатору.
\end{itemize}

Зависимости:
\begin{itemize}
	\item взаимодействует с сущностями, определенными в пространстве имен App/Entity/Vacancy/Vacancy, через ManagerRegistry для динамического получения репозитория;
	\item использует FOSRestBundle для обработки REST-запросов и формирования JSON-ответов;
	\item применяет Symfony/Component/HttpKernel/Exception/NotFoundHttpException для обработки ошибок при неверном типе параметра.
\end{itemize}

\begin{xltabular}{\textwidth}{|X|X|}
	\caption{Данные класса VacancyParametersApiController \label{VacancyPar:table}} \\
	\hline 
	\centrow Данные & 
	\centrow Описание \\
	\hline 
	\endfirsthead
	
	\caption*{Продолжение таблицы \ref{VacancyPar:table}} \\
	\hline 
	\centrow Данные & 
	\centrow Описание  \\
	\hline 
	\endhead
	
	vacancyParameter &  Строка, определяющая тип параметра вакансии (часть имени класса сущности). \\ \hline
	ManagerRegistry &  Интерфейс для получения репозитория сущностей Doctrine. \\ \hline
	NotFoundHttpException &  Исключение, выбрасываемое при отсутствии запрошенного типа параметра. \\ \hline
\end{xltabular}

\begin{xltabular}{\textwidth}{|X|X|}
	\caption{Методы класса  VacancyParametersApiController \label{VacancypM:table}} \\
	\hline 
	\centrow Метод & 
	\centrow Описание \\
	\hline 
	\endfirsthead
	
	\caption*{Продолжение таблицы \ref{VacancypM:table}} \\
	\hline 
	\centrow Метод & 
	\centrow Описание  \\
	\hline 
	\endhead
	
	getManyVacancyParameters(string \$vacancyParameter, ManagerRegistry \$doctrine) &  Возвращает список всех параметров вакансии для указанного типа. \\ \hline
	getSingleVacancyParameters(string \$vacancyParameter, int \$id, ManagerRegistry \$doctrine) &  Возвращает данные об одном параметре вакансии по его идентификатору. \\ \hline
\end{xltabular}

\subsubsection{Спецификация класса ResumeApiController}
ResumeApiController (компонент управления резюме через API)

Описание: Управляет операциями, связанными с резюме, включая создание, получение списка резюме, получение данных об одном резюме и получение резюме, связанных с авторизованным пользователем. Реализует REST API и взаимодействует с сущностью Resume.

Валидация данных:
\begin{itemize}
	\item в методе newResume() используется форма ResumeApiType для валидации входных данных, полученных в формате JSON;
	\item проверяет корректность полей резюме перед сохранением в базе данных.
\end{itemize}

Связи:
\begin{itemize}
	\item резюме связано с сущностями WorkPlace и Education, которые привязываются к резюме в методе newResume();
	\item резюме связано с пользователем через поле user, устанавливаемое автоматически для авторизованного пользователя.
\end{itemize}

Безопасность:
\begin{itemize}
	\item использует FOSRestBundle для обработки REST-запросов и формирования стандартизированных JSON-ответов;
	\item доступ к методу userResume() ограничен для авторизованных пользователей через getUser().
\end{itemize}

Основные методы:
\begin{itemize}
	\item getManyResumes(ResumeRepository \$resumeRepository, Request \$request): возвращает отфильтрованный список резюме с пагинацией и информацией о фильтрах;
	\item getSingleResume(Resume \$resume, ResumeRepository \$resumeRepository): возвращает данные об одном резюме по его идентификатору;
	\item newResume(Request \$request, EntityManagerInterface \$entityManager): создает новое резюме, валидирует данные и сохраняет их в базе данных;
	\item userResume(ResumeRepository \$resumeRepository): возвращает список резюме, связанных с авторизованным пользователем;
\end{itemize}

Зависимости:
\begin{itemize}
	\item взаимодействует с сущностью Resume и репозиторием ResumeRepository;
	\item использует ResumeApiType для обработки форм создания резюме;
	\item зависит от Doctrine/ORM/EntityManagerInterface для работы с базой данных;
	\item применяет FilterService для обработки и применения фильтров в методе getManyResumes;
	\item использует FOSRestBundle для обработки REST-запросов и формирования JSON-ответов;
\end{itemize}

\begin{xltabular}{\textwidth}{|X|X|}
	\caption{Данные класса ResumeApiController \label{Resume:table}} \\
	\hline 
	\centrow Данные & 
	\centrow Описание \\
	\hline 
	\endfirsthead
	
	\caption*{Продолжение таблицы \ref{Resume:table}} \\
	\hline 
	\centrow Данные & 
	\centrow Описание  \\
	\hline 
	\endhead
	
	Request &  Данные из HTTP-запросов (JSON) для валидации и обработки форм. \\ \hline
	Resume &  Сущность, представляющая резюме. \\ \hline
	ResumeApiType &  Форма для валидации данных при создании резюме. \\ \hline
	FilterService &  Сервис для обработки и применения фильтров к списку резюме. \\ \hline
	EntityManagerInterface &  Интерфейс для работы с базой данных через Doctrine. \\ \hline
\end{xltabular}

\begin{xltabular}{\textwidth}{|X|X|}
	\caption{Методы класса  VacancyParametersApiController \label{ResumeM:table}} \\
	\hline 
	\centrow Метод & 
	\centrow Описание \\
	\hline 
	\endfirsthead
	
	\caption*{Продолжение таблицы \ref{ResumeM:table}} \\
	\hline 
	\centrow Метод & 
	\centrow Описание  \\
	\hline 
	\endhead
	
	getManyResumes(ResumeRepository \$resumeRepository, Request \$request) &  Возвращает отфильтрованный список резюме с пагинацией и информацией о фильтрах. \\ \hline
	getSingleResume(Resume \$resume, ResumeRepository \$resumeRepository) &  Возвращает данные об одном резюме по его идентификатору. \\ \hline
	newResume(Request \$request, EntityManagerInterface \$entityManager) &  Валидирует входные данные, сохраняет новое резюме в базе данных. \\ \hline
	userResume(ResumeRepository \$resumeRepository) &  Возвращает список резюме, связанных с авторизованным пользователем. \\ \hline
\end{xltabular}

\subsubsection{Спецификация класса ResumeParametersApiController}
ResumeParametersApiController (компонент управления параметрами резюме через API)

Описание: Управляет операциями получения параметров резюме через REST API. Обеспечивает получение списка всех параметров определенного типа или данных об одном параметре по его идентификатору. Взаимодействует с сущностями, связанными с параметрами резюме.

Валидация данных:
\begin{itemize}
	\item отсутствует валидация входных данных, так как методы обрабатывают GET-запросы с параметрами, передаваемыми в URL;
	\item проверяется существование запрошенного типа параметра резюме, в случае ошибки выбрасывается исключение NotFoundHttpException.
\end{itemize}

Связи:
\begin{itemize}
	\item параметры резюме связаны с сущностями в пространстве имен App/Entity/Resume, определяемыми динамически через параметр resumeParameter;
\end{itemize}

Безопасность:
\begin{itemize}
	\item использует FOSRestBundle для обработки REST-запросов и формирования стандартизированных JSON-ответов;
	\item не реализует дополнительной аутентификации или авторизации, доступ к методам предполагается открытым.
\end{itemize}


Основные методы:
\begin{itemize}
	\item getManyResumeParameters(string \$resumeParameter, ManagerRegistry \$doctrine): возвращает список всех параметров резюме для указанного типа;
	\item getSingleResumeParameters(string \$resumeParameter, int \$id, ManagerRegistry \$doctrine): возвращает данные об одном параметре резюме по его идентификатору.
\end{itemize}

Зависимости:
\begin{itemize}
	\item взаимодействует с сущностями, определенными в пространстве имен App/Entity/Resume, через ManagerRegistry для динамического получения репозитория;
	\item использует FOSRestBundle для обработки REST-запросов и формирования JSON-ответов;
	\item применяет Symfony/Component/HttpKernel/Exception/NotFoundHttpException для обработки ошибок при неверном типе параметра.
\end{itemize}

\begin{xltabular}{\textwidth}{|X|X|}
	\caption{Данные класса ResumeParametersApiController \label{ResumePar:table}} \\
	\hline 
	\centrow Данные & 
	\centrow Описание \\
	\hline 
	\endfirsthead
	
	\caption*{Продолжение таблицы \ref{ResumePar:table}} \\
	\hline 
	\centrow Данные & 
	\centrow Описание  \\
	\hline 
	\endhead
	
	resumeParameter & Строка, определяющая тип параметра резюме (часть имени класса сущности). \\ \hline
	ManagerRegistry & Интерфейс для получения репозитория сущностей Doctrine. \\ \hline
	NotFoundHttpException & Исключение, выбрасываемое при отсутствии запрошенного типа параметра. \\ \hline
\end{xltabular}

\begin{xltabular}{\textwidth}{|X|X|}
	\caption{Методы класса ResumeParametersApiController \label{ResumePM:table}} \\
	\hline 
	\centrow Метод & 
	\centrow Описание \\
	\hline 
	\endfirsthead
	
	\caption*{Продолжение таблицы \ref{ResumePM:table}} \\
	\hline 
	\centrow Метод & 
	\centrow Описание  \\
	\hline 
	\endhead
	
	getManyResumeParameters(string \$resumeParameter, ManagerRegistry \$doctrine) & Возвращает список всех параметров резюме для указанного типа. \\ \hline
	getSingleResumeParameters(string \$resumeParameter, int \$id, ManagerRegistry \$doctrine) & Возвращает данные об одном параметре резюме по его идентификатору. \\ \hline
\end{xltabular}

\subsubsection{Спецификация класса CompanyApiController}
CompanyApiController (компонент управления компаниями через API)

Описание: Управляет операциями, связанными с компаниями, включая создание, получение списка компаний, получение данных об одной компании, обновление и удаление компаний, а также получение компаний, связанных с авторизованным пользователем. Реализует REST API и взаимодействует с сущностью Company.

Валидация данных:
\begin{itemize}
	\item в методах newCompany() и updateCompany() используется форма CompanyApiType для валидации входных данных, полученных в формате JSON;
	\item проверяет корректность полей компании перед сохранением в базе данных.
\end{itemize}

Связи:
\begin{itemize}
	\item компании связаны с пользователем через поле user, устанавливаемое автоматически для авторизованного пользователя.
\end{itemize}

Безопасность:
\begin{itemize}
	\item использует FOSRestBundle для обработки REST-запросов и формирования стандартизированных JSON-ответов;
	\item проверяет роль ROLE\_EMPLOYER для создания компании в методе newCompany();
	\item ограничивает доступ к методам updateCompany() и deleteCompany(), проверяя, является ли пользователь владельцем компании;
	\item доступ к методу userCompanies() ограничен для авторизованных пользователей через getUser().
\end{itemize}

Основные методы:
\begin{itemize}
	\item getManyCompanies(CompanyRepository \$companyRepository, Request \$request): возвращает отфильтрованный список компаний с пагинацией и информацией о фильтрах;
	\item getSingleCompany(Company \$company, CompanyRepository \$companyRepository): возвращает данные об одной компании по её идентификатору;
	\item newCompany(Request \$request, EntityManagerInterface \$entityManager): создает новую компанию, валидирует данные и сохраняет их в базе данных;
	\item updateCompany(Company \$company, Request \$request, EntityManagerInterface \$entityManager): обновляет данные компании после валидации;
	\item deleteCompany(Company \$company, EntityManagerInterface \$entityManager): удаляет компанию из базы данных;
	\item userCompanies(CompanyRepository \$companyRepository): возвращает список компаний, связанных с авторизованным пользователем.
\end{itemize}

Зависимости:
\begin{itemize}
	\item взаимодействует с сущностью Company и репозиторием CompanyRepository;
	\item использует CompanyApiType для обработки форм создания и обновления компаний;
	\item зависит от Doctrine/ORM/EntityManagerInterface для работы с базой данных;
	\item применяет FilterService для обработки и применения фильтров в методе getManyCompanies;
	\item использует FOSRestBundle для обработки REST-запросов и формирования JSON-ответов;
	\item применяет Symfony/Component/HttpKernel/Exception/AccessDeniedHttpException для обработки ошибок доступа.
\end{itemize}

\begin{xltabular}{\textwidth}{|X|X|}
	\caption{Данные класса CompanyApiController \label{CompanyPar:table}} \\
	\hline 
	\centrow Данные & 
	\centrow Описание \\
	\hline 
	\endfirsthead
	
	\caption*{Продолжение таблицы \ref{CompanyPar:table}} \\
	\hline 
	\centrow Данные & 
	\centrow Описание  \\
	\hline 
	\endhead
	
	Request & Данные из HTTP-запросов (JSON) для валидации и обработки форм. \\ \hline
	Company & Сущность, представляющая компанию. \\ \hline
	CompanyApiType & Форма для валидации данных при создании и обновлении компании. \\ \hline
	FilterService & Сервис для обработки и применения фильтров к списку компаний. \\ \hline
	EntityManagerInterface & Интерфейс для работы с базой данных через Doctrine. \\ \hline
	AccessDeniedHttpException & Исключение, выбрасываемое при отсутствии прав доступа. \\ \hline
\end{xltabular}

\begin{xltabular}{\textwidth}{|X|X|}
	\caption{Методы класса CompanyApiController \label{CompanyPM:table}} \\
	\hline 
	\centrow Метод & 
	\centrow Описание \\
	\hline 
	\endfirsthead
	
	\caption*{Продолжение таблицы \ref{CompanyPM:table}} \\
	\hline 
	\centrow Метод & 
	\centrow Описание  \\
	\hline 
	\endhead
	
	getManyCompanies(CompanyRepository \$companyRepository, Request \$request) & Возвращает отфильтрованный список компаний с пагинацией и информацией о фильтрах в формате JSON с кодом 200. \\ \hline
	getSingleCompany(Company \$company, CompanyRepository \$companyRepository) & Возвращает данные об одной компании по её идентификатору в формате JSON с кодом 200. \\ \hline
	newCompany(Request \$request, EntityManagerInterface \$entityManager) & Валидирует входные данные, сохраняет новую компанию в базе данных. Возвращает HTTP-ответ с кодом 201 при успехе или 400 при ошибке. \\ \hline
	updateCompany(Company \$company, Request \$request, EntityManagerInterface \$entityManager) & Обновляет данные компании после валидации. Возвращает HTTP-ответ с кодом 200 при успехе, 400 при ошибке валидации или 403 при отсутствии прав. \\ \hline
	deleteCompany(Company \$company, EntityManagerInterface \$entityManager) & Удаляет компанию из базы данных. Возвращает HTTP-ответ с кодом 200 при успехе или 403 при отсутствии прав. \\ \hline
	userCompanies(CompanyRepository \$companyRepository) & Возвращает список компаний, связанных с авторизованным пользователем, в формате JSON с кодом 200. \\ \hline
\end{xltabular}

\subsubsection{Спецификация класса VacancyResponseApiController}
VacancyResponseApiController (компонент управления откликами на вакансии через API)

Описание: Управляет операциями, связанными с откликами на вакансии, включая создание нового отклика и получение списка откликов, связанных с авторизованным пользователем. Реализует REST API и взаимодействует с сущностью VacancyResponse.

Валидация данных:
\begin{itemize}
	\item в методе newVacancyResponse() используется форма VacancyResponseApiType для валидации входных данных, полученных в формате JSON;
	\item проверяет наличие активных резюме у пользователя перед созданием отклика;
	\item при отсутствии активных резюме выбрасывается исключение HttpException с кодом 300.
\end{itemize}

Связи:
\begin{itemize}
	\item отклики на вакансии связаны с сущностями Resume, VacancyResponseStatus и пользователем через поле initiator.
\end{itemize}

Безопасность:
\begin{itemize}
	\item использует FOSRestBundle для обработки REST-запросов и формирования стандартизированных JSON-ответов;
	\item доступ к методу userResume() ограничен для авторизованных пользователей через getUser();
	\item метод newVacancyResponse() доступен только пользователям с активными резюме.
\end{itemize}

Уведомления:
\begin{itemize}
	\item не реализует уведомления для пользователей, так как методы возвращают только HTTP-ответы в формате JSON;
	\item в случае отсутствия активных резюме возвращается сообщение об ошибке в формате JSON.
\end{itemize}

Основные методы:
\begin{itemize}
	\item newVacancyResponse(Request \$request, VacancyResponseStatusRepository \$vacancyResponseStatusRepository, EntityManagerInterface \$entityManager, UserRepository \$userRepository): создает новый отклик на вакансию, валидирует данные и сохраняет их в базе данных;
	\item userResume(VacancyResponseRepository \$vacancyResponseRepository): возвращает список откликов, связанных с авторизованным пользователем.
\end{itemize}

Зависимости:
\begin{itemize}
	\item взаимодействует с сущностями VacancyResponse, Resume, VacancyResponseStatus и репозиториями VacancyResponseRepository, VacancyResponseStatusRepository, UserRepository;
	\item использует VacancyResponseApiType для обработки форм создания откликов;
	\item зависит от Doctrine/ORM/EntityManagerInterface для работы с базой данных;
	\item использует FOSRestBundle для обработки REST-запросов и формирования JSON-ответов;
	\item применяет Symfony/Component/HttpKernel/Exception/HttpException для обработки ошибок при отсутствии активных резюме.
\end{itemize}

\begin{xltabular}{\textwidth}{|X|X|}
	\caption{Данные класса VacancyResponseApiController \label{VacancyResponsePar:table}} \\
	\hline 
	\centrow Данные & 
	\centrow Описание \\
	\hline 
	\endfirsthead
	
	\caption*{Продолжение таблицы \ref{VacancyResponsePar:table}} \\
	\hline 
	\centrow Данные & 
	\centrow Описание  \\
	\hline 
	\endhead
	
	Request & Данные из HTTP-запросов (JSON) для валидации и обработки форм. \\ \hline
	VacancyResponse & Сущность, представляющая отклик на вакансию. \\ \hline
	Resume & Сущность, представляющая резюме пользователя. \\ \hline
	VacancyResponseStatus & Сущность, представляющая статус отклика на вакансию. \\ \hline
	VacancyResponseApiType & Форма для валидации данных при создании отклика. \\ \hline
	EntityManagerInterface & Интерфейс для работы с базой данных через Doctrine. \\ \hline
	HttpException & Исключение, выбрасываемое при отсутствии активных резюме. \\ \hline
\end{xltabular}

\begin{xltabular}{\textwidth}{|X|X|}
	\caption{Методы класса VacancyResponseApiController \label{VacancyResponsePM:table}} \\
	\hline 
	\centrow Метод & 
	\centrow Описание \\
	\hline 
	\endfirsthead
	
	\caption*{Продолжение таблицы \ref{VacancyResponsePM:table}} \\
	\hline 
	\centrow Метод & 
	\centrow Описание  \\
	\hline 
	\endhead
	
	newVacancyResponse(Request \$request, VacancyResponseStatusRepository \$vacancyResponseStatusRepository, EntityManagerInterface \$entityManager, UserRepository \$userRepository) & Валидирует входные данные, сохраняет новый отклик на вакансию в базе данных. Возвращает HTTP-ответ с кодом 201 при успехе, 400 при ошибке валидации или 300 при отсутствии активных резюме. \\ \hline
	userResume(VacancyResponse Repository \$vacancyResponseRepository) & Возвращает список откликов, связанных с авторизованным пользователем, в формате JSON с кодом 200. \\ \hline
\end{xltabular}

\subsubsection{Спецификация класса CityApiController}
CityApiController (компонент управления городами через API)

Описание: Управляет операциями получения информации о городах через REST API. Обеспечивает получение списка всех городов или данных об одном городе по его идентификатору. Взаимодействует с сущностью City.

Валидация данных:
\begin{itemize}
	\item отсутствует валидация входных данных, так как методы обрабатывают GET-запросы с параметрами, передаваемыми в URL;
	\item предполагается, что идентификатор города, переданный в URL, существует в базе данных.
\end{itemize}


Безопасность:
\begin{itemize}
	\item использует FOSRestBundle для обработки REST-запросов и формирования стандартизированных JSON-ответов;
	\item не реализует дополнительной аутентификации или авторизации, доступ к методам предполагается открытым.
\end{itemize}


Основные методы:
\begin{itemize}
	\item getCityList(CityRepository \$cityRepository): возвращает список всех городов;
	\item getSingleCity(City \$city, CityRepository \$cityRepository): возвращает данные об одном городе по его идентификатору.
\end{itemize}

Зависимости:
\begin{itemize}
	\item взаимодействует с сущностью City и репозиторием CityRepository;
	\item использует FOSRestBundle для обработки REST-запросов и формирования JSON-ответов.
\end{itemize}

\begin{xltabular}{\textwidth}{|X|X|}
	\caption{Данные класса CityApiController \label{CityPar:table}} \\
	\hline 
	\centrow Данные & 
	\centrow Описание \\
	\hline 
	\endfirsthead
	
	\caption*{Продолжение таблицы \ref{CityPar:table}} \\
	\hline 
	\centrow Данные & 
	\centrow Описание  \\
	\hline 
	\endhead
	
	City & Сущность, представляющая город. \\ \hline
	CityRepository & Репозиторий для работы с сущностью City. \\ \hline
\end{xltabular}

\begin{xltabular}{\textwidth}{|X|X|}
	\caption{Методы класса CityApiController \label{CityPM:table}} \\
	\hline 
	\centrow Метод & 
	\centrow Описание \\
	\hline 
	\endfirsthead
	
	\caption*{Продолжение таблицы \ref{CityPM:table}} \\
	\hline 
	\centrow Метод & 
	\centrow Описание  \\
	\hline 
	\endhead
	
	getCityList(CityRepository \$cityRepository) & Возвращает список всех городов в формате JSON с кодом 200. \\ \hline
	getSingleCity(City \$city, CityRepository \$cityRepository) & Возвращает данные об одном городе по его идентификатору в формате JSON с кодом 200. \\ \hline
\end{xltabular}

\subsubsection{Спецификация класса DoctrineFiltersEnabledSubscriber}
DoctrineFiltersEnabledSubscriber (сервис для управления фильтрами Doctrine)

Описание: Отвечает за динамическое включение фильтров Doctrine для ограничения данных, возвращаемых из базы данных, на основе маршрутов и консольных команд. Реализует интерфейс EventSubscriberInterface для обработки событий HTTP-запросов. Инкапсулирует логику управления фильтрами, связанными с активностью и модерацией сущностей.

Валидация данных:
\begin{itemize}
	\item отсутствует валидация входных данных, так как методы обрабатывают параметры маршрутов и URI;
	\item проверяет наличие маршрутов в белом списке и присутствие префикса /admin в URI.
\end{itemize}

Связи:
\begin{itemize}
	\item взаимодействует с EntityManagerInterface для управления фильтрами Doctrine.
\end{itemize}

Безопасность:
\begin{itemize}
	\item применяет фильтры only\_enabled\_from\_front\_doctrine\_filter и only\_moderated\_from\_front\_doctrine\_filter для ограничения доступа к данным на уровне базы данных;
	\item исключает применение фильтров для маршрутов, указанных в белых списках, и запросов с префиксом /admin.
\end{itemize}

Основные методы:
\begin{itemize}
	\item getSubscribedEvents(): возвращает массив подписанных событий (KernelEvents::REQUEST и ConsoleEvents::COMMAND);
	\item onKernelRequest(RequestEvent \$event): обрабатывает HTTP-запросы, вызывая методы проверки и включения фильтров;
	\item onConsoleCommand(ConsoleCommandEvent \$event): активирует фильтры для определенной консольной команды (app:search-import);
	\item checkAndEnabledFilterIsEnabled(string \$uri, ?string \$routeNamed): проверяет условия и включает фильтр is\_enabled;
	\item checkAndEnabledFilterIsModerated(string \$uri, ?string \$routeNamed): проверяет условия и включает фильтр is\_moderated;
	\item setEnabledFilterIsEnabled(): активирует фильтр only\_enabled\_from\_front\_doctrine\_filter с параметром is\_enabled;
	\item setEnabledFilterIsModerated(): активирует фильтр only\_moderated\_from\_front\_doctrine\_filter с параметром is\_moderated.
\end{itemize}

Зависимости:
\begin{itemize}
	\item использует EntityManagerInterface для работы с фильтрами Doctrine;
	\item зависит от Symfony/Component/HttpKernel/Event/RequestEvent и Symfony/Component/Console/Event/ConsoleCommandEvent для обработки событий;
	\item взаимодействует с FOSRestBundle для обработки REST-запросов в контексте маршрутов.
\end{itemize}

\begin{xltabular}{\textwidth}{|X|X|}
	\caption{Данные класса DoctrineFiltersEnabledSubscriber \label{DoctrineFiltersPar:table}} \\
	\hline 
	\centrow Данные & 
	\centrow Описание \\
	\hline 
	\endfirsthead
	
	\caption*{Продолжение таблицы \ref{DoctrineFiltersPar:table}} \\
	\hline 
	\centrow Данные & 
	\centrow Описание  \\
	\hline 
	\endhead
	
	EntityManagerInterface & Интерфейс для управления фильтрами Doctrine. \\ \hline
	RequestEvent & Событие HTTP-запроса для обработки маршрутов и URI. \\ \hline
	ConsoleCommandEvent & Событие консольной команды для обработки команды app:search-import. \\ \hline
\end{xltabular}

\begin{xltabular}{\textwidth}{|X|X|}
	\caption{Методы класса DoctrineFiltersEnabledSubscriber \label{DoctrineFiltersPM:table}} \\
	\hline 
	\centrow Метод & 
	\centrow Описание \\
	\hline 
	\endfirsthead
	
	\caption*{Продолжение таблицы \ref{DoctrineFiltersPM:table}} \\
	\hline 
	\centrow Метод & 
	\centrow Описание  \\
	\hline 
	\endhead
	
	getSubscribedEvents() & Возвращает массив подписанных событий для обработки HTTP-запросов и консольных команд. \\ \hline
	onKernelRequest(RequestEvent \$event) & Обрабатывает HTTP-запросы, вызывая методы проверки и включения фильтров. \\ \hline
	onConsoleCommand(Console
	CommandEvent \$event) & Активирует фильтры для консольной команды app:search-import. \\ \hline
	checkAndEnabledFilterIsEnabled( string \$uri, ?string \$routeNamed) & Проверяет условия и включает фильтр is\_enabled для ограничения данных. \\ \hline
	checkAndEnabledFilterIsModerated( string \$uri, ?string \$routeNamed) & Проверяет условия и включает фильтр is\_moderated для ограничения данных. \\ \hline
	setEnabledFilterIsEnabled() & Активирует фильтр only\_enabled\_from\_front\_doctrine \_filter с параметром is\_enabled. \\ \hline
	setEnabledFilterIsModerated() & Активирует фильтр only\_moderated\_from\_front\_doctrine \_filter с параметром is\_moderated. \\ \hline
\end{xltabular}

\subsubsection{Спецификация класса FilterService}
FilterService (сервис для управления фильтрацией данных)

Описание: Отвечает за обработку и применение фильтров к запросам Doctrine на основе параметров HTTP-запросов. Парсит фильтры из запросов, определяет фильтруемые поля сущностей с помощью атрибута Filterable и применяет фильтры к QueryBuilder. Инкапсулирует логику фильтрации данных для выборки сущностей из базы данных.

Основные задачи:
\begin{itemize}
	\item парсинг параметров фильтров из HTTP-запросов;
	\item определение фильтруемых полей сущностей через атрибут Filterable;
	\item применение фильтров к запросам Doctrine с поддержкой различных типов фильтрации (точное совпадение, LIKE, диапазоны, даты);
	\item обработка связанных сущностей через LEFT JOIN;
	\item проверка валидности значений фильтров и их преобразование (например, диапазоны вида "100-500").
\end{itemize}

Основные методы:
\begin{itemize}
	\item parseFilters(Request \$request): парсит параметры фильтров из HTTP-запроса и возвращает их в виде массива;
	\item applyFilters(QueryBuilder \$qb, array \$filters, string \$entityClass): применяет фильтры к QueryBuilder на основе фильтруемых полей сущности;
	\item getFilterableFields(ReflectionClass \$reflection): возвращает массив фильтруемых полей сущности, определенных атрибутом Filterable;
	\item applyFilter(QueryBuilder \$qb, string \$type, string \$field, mixed \$value, string \$paramName): применяет фильтр к QueryBuilder в зависимости от типа фильтра;
	\item applyRangeFilter(QueryBuilder \$qb, string \$field, mixed \$value, string \$paramName): применяет фильтр диапазона значений;
	\item applyDateRangeFilter(QueryBuilder \$qb, string \$field, mixed \$value, string \$paramName): применяет фильтр диапазона дат;
	\item addJoin(QueryBuilder \$qb, string \$joinConfig): добавляет LEFT JOIN к QueryBuilder для связанных сущностей;
	\item getFieldPath(Filterable \$filterable, string \$rootAlias): определяет путь к полю для фильтрации;
	\item isValidFilterValue(mixed \$value): проверяет валидность значения фильтра;
	\item parseFilterValue(mixed \$value): парсит значение фильтра, включая диапазоны.
\end{itemize}

Зависимости:
\begin{itemize}
	\item использует EntityManagerInterface для работы с Doctrine и QueryBuilder;
	\item зависит от Symfony/Component/HttpFoundation/Request для получения параметров запроса;
	\item использует ReflectionClass для анализа полей сущностей;
	\item взаимодействует с атрибутом App/Attribute/Filterable для определения фильтруемых полей;
	\item использует App/Enum/FilterType для определения типов фильтров.
\end{itemize}

\begin{xltabular}{\textwidth}{|X|X|}
	\caption{Данные класса FilterService \label{FilterServicePar:table}} \\
	\hline 
	\centrow Данные & 
	\centrow Описание \\
	\hline 
	\endfirsthead
	
	\caption*{Продолжение таблицы \ref{FilterServicePar:table}} \\
	\hline 
	\centrow Данные & 
	\centrow Описание  \\
	\hline 
	\endhead
	
	EntityManagerInterface & Интерфейс для работы с Doctrine и QueryBuilder. \\ \hline
	Request & Объект HTTP-запроса для получения параметров фильтров. \\ \hline
	ReflectionClass & Класс для анализа полей сущностей через рефлексию. \\ \hline
	Filterable & Атрибут, определяющий фильтруемые поля сущности. \\ \hline
	FilterType & Перечисление, определяющее типы фильтров. \\ \hline
\end{xltabular}

\begin{xltabular}{\textwidth}{|X|X|}
	\caption{Методы класса FilterService \label{FilterServicePM:table}} \\
	\hline 
	\centrow Метод & 
	\centrow Описание \\
	\hline 
	\endfirsthead
	
	\caption*{Продолжение таблицы \ref{FilterServicePM:table}} \\
	\hline 
	\centrow Метод & 
	\centrow Описание  \\
	\hline 
	\endhead
	
	parseFilters(Request \$request) & Парсит параметры фильтров из HTTP-запроса и возвращает их в виде массива. \\ \hline
	applyFilters(QueryBuilder \$qb, array \$filters, string \$entityClass) & Применяет фильтры к QueryBuilder на основе фильтруемых полей сущности. \\ \hline
	getFilterableFields(ReflectionClass \$reflection) & Возвращает массив фильтруемых полей сущности, определенных атрибутом Filterable. \\ \hline
	applyFilter(QueryBuilder \$qb, string \$type, string \$field, mixed \$value, string \$paramName) & Применяет фильтр к QueryBuilder в зависимости от типа фильтра. \\ \hline
	applyRangeFilter(QueryBuilder \$qb, string \$field, mixed \$value, string \$paramName) & Применяет фильтр диапазона значений. \\ \hline
	applyDateRangeFilter(QueryBuilder \$qb, string \$field, mixed \$value, string \$paramName) & Применяет фильтр диапазона дат. \\ \hline
	addJoin(QueryBuilder \$qb, string \$joinConfig) & Добавляет LEFT JOIN к QueryBuilder для связанных сущностей. \\ \hline
	getFieldPath(Filterable \$filterable, string \$rootAlias) & Определяет путь к полю для фильтрации. \\ \hline
	isValidFilterValue(mixed \$value) & Проверяет валидность значения фильтра. \\ \hline
	parseFilterValue(mixed \$value) & Парсит значение фильтра, включая диапазоны. \\ \hline
\end{xltabular}

\subsubsection{Спецификация класса StringHelperService}
StringHelperService (Сервис для обработки строк)

Описание: Отвечает за обработку и преобразование строк, включая генерацию случайных строк, создание слаг-идентификаторов, транслитерацию русского текста в английский, извлечение аргументов из шаблонов строк и преобразование регистра. Инкапсулирует вспомогательные функции для работы со строками в приложении.

Основные задачи:
\begin{itemize}
	\item генерация случайных строк заданной длины;
	\item создание уникальных слаг-идентификаторов для сущностей;
	\item транслитерация русских символов в латинские для создания читаемых URL;
	\item извлечение аргументов из шаблонов строк формата \%argument\$s;
	\item преобразование первой буквы строки в верхний регистр с учетом многобайтных кодировок.
\end{itemize}

Основные методы:
\begin{itemize}
	\item getRandomStr(int \$length): генерирует случайную строку заданной длины из набора символов;
	\item generateSlugForEntity(string \$string, bool \$unique = true): создает слаг из строки с возможным добавлением уникального суффикса;
	\item translitRuToEn(string \$text): транслитерирует русский текст в английский, заменяя специальные символы на дефисы;
	\item getStringArguments(\$format): извлекает аргументы из шаблона строки формата \%argument\$s;
	\item my\_mb\_ucfirst(string \$string): преобразует первую букву строки в верхний регистр с учетом многобайтных кодировок.
\end{itemize}

Зависимости:
\begin{itemize}
	\item использует встроенные PHP-функции для работы со строками (mb\_strtolower, mb\_ereg\_replace, strtr);
	\item зависит от Exception для обработки ошибок при генерации случайных строк;
	\item не имеет прямых зависимостей от внешних библиотек или фреймворков.
\end{itemize}

\begin{xltabular}{\textwidth}{|X|X|}
	\caption{Данные класса StringHelperService \label{StringHelperServicePar:table}} \\
	\hline 
	\centrow Данные & 
	\centrow Описание \\
	\hline 
	\endfirsthead
	
	\caption*{Продолжение таблицы \ref{StringHelperServicePar:table}} \\
	\hline 
	\centrow Данные & 
	\centrow Описание  \\
	\hline 
	\endhead
	
	KEYSPACE & Строка, содержащая набор символов для генерации случайных строк. \\ \hline
\end{xltabular}

\begin{xltabular}{\textwidth}{|X|X|}
	\caption{Методы класса StringHelperService \label{StringHelperServicePM:table}} \\
	\hline 
	\centrow Метод & 
	\centrow Описание \\
	\hline 
	\endfirsthead
	
	\caption*{Продолжение таблицы \ref{StringHelperServicePM:table}} \\
	\hline 
	\centrow Метод & 
	\centrow Описание  \\
	\hline 
	\endhead
	
	getRandomStr(int \$length) & Генерирует случайную строку заданной длины из набора символов. \\ \hline
	generateSlugForEntity(string \$string, bool \$unique = true) & Создает слаг из строки с возможным добавлением уникального суффикса. \\ \hline
	translitRuToEn(string \$text) & Транслитерирует русский текст в английский, заменяя специальные символы на дефисы. \\ \hline
	getStringArguments(\$format) & Извлекает аргументы из шаблона строки формата \%argument\$s. \\ \hline
	my\_mb\_ucfirst(string \$string) & Преобразует первую букву строки в верхний регистр с учетом многобайтных кодировок. \\ \hline
\end{xltabular}


















\subsection{Тестирование программной системы}

\subsubsection{Модульное тестирование классов и методов программы}

Модульное тестирование в рамках разработки программноинформаци-
онной системы для управления данными о подготовке к производству метал-
локонструкций проводится с целью выявления ошибок в работе отдельных
компонентов программы. Оно позволяет убедиться, что каждый класс и ме-
тод системы функционируют корректно при различных сценариях использо-
вания.

Описание процесса модульного тестирования:

Модульное тестирование фокусируется на проверке правильности ра-
боты небольших единиц программы — классов и методов, что помогает
быстро локализовать и устранять ошибки на ранних этапах разработки. В
данной системе тестирование охватывает такие ключевые компоненты, как
управление пользователями, обработка заказов, работа с чертежами, матери-
алами, задачами и отчетами, а также взаимодействие с хранилищем данных
и другими внешними компонентами.

Основные цели модульного тестирования:
\begin{enumerate}
	\item Изоляция: Каждый тест проверяет работу одного класса или метода,
	что позволяет точно определить место возникновения ошибки.
	\item Автоматизация: Модульные тесты выполняются автоматически с
	использованием таких инструментов, как PHPUnit (для PHP), что ускоряет
	процесс разработки и исправления ошибок.
	\item Покрытие: Тесты охватывают все ключевые методы системы, вклю-
	чая проверку обработки граничных случаев (например, пустые или некор-
	ректные входные данные).
\end{enumerate}
Модульное тестирование для различных классов системы:

Тестирование класса UserApiController 

\begin{enumerate}
	\item Тест метода создания пользователя(‘userRegistration‘):
	\begin{itemize}
		\item проверить, что при передаче корректных данных новый пользователь
		сохраняется в базе данных;
		\item проверить, что при отсутствии обязательных данных (например, пу-
		стое имя или email) возвращается ошибка валидации.
		
		\begin{figure}[H]
			\begin{lstlisting}[language=PHP]
				public function testRegistration()
				{
					$data = [
					'firstName' => 'testFirstName',
					'lastName' => 'testLastName',
					'password' => 'testPassword',
					'email' => 'testa@test.com',
					'phone' => '89192166761'
					];
					
					$response = $this->apiRequest(
					$this->router->generate('api_registration'),
					'POST',
					content: json_encode($data),
					);
					
					$this->assertEquals(201, $response->getStatusCode());
					$userData = $this->em->getRepository(User::class)->findOneBy(['email' => 'testa@test.com']);
					$this->assertEquals($data['firstName'], $userData->getFirstName());
					$this->assertEquals($data['lastName'], $userData->getLastName());
					$this->assertEquals($data['email'], $userData->getEmail());
				}
				
				public function testRegistrationWrong()
				{
					$data = [
					'firstName' => '',
					'lastName' => '',
					'password' => '',
					'email' => '',
					'phone' => ''
					];
					
					$response = $this->apiRequest(
					$this->router->generate('api_registration'),
					'POST',
					content: json_encode($data),
					);
					
					$this->assertEquals(400, $response->getStatusCode());
					$content = json_decode($response->getContent(), true);
					$this->assertEquals('Validation Failed', $content['form']['message']);
				}
			\end{lstlisting}  
			\caption{Модульное тестирование метода создания пользователя
				класса UserApiController}
			\label{unitUser:image}
		\end{figure}
	\end{itemize}
\end{enumerate}





Тестирование класса VacancyApiController 

\begin{enumerate}
	\item Тест метода создания вакансии(‘newVacancy‘):
	\begin{itemize}
		\item проверить, что при передаче корректных данных вакансия
		сохраняется в базе данных;
		\item проверить, что при отсутствии обязательных данных возвращается ошибка валидации.
		\begin{figure}[H]
			\begin{lstlisting}[language=PHP]
				public function testApiNewVacancy()
				{
					$this->createUser();
					$this->em->flush();
					$data = [
					'name' => 'test', 'workAddress' => 'test', 'description' => 'test',
					'specializations' => [1],
					'isEnabled' => true,
					];
					$response = $this->apiRequest(
					$this->router->generate('api_create_new_vacancy'),
					'POST',
					headers: [
					'HTTP_AUTHORIZATION' => sprintf(
					'Bearer %s',
					$this->getAuthToken(self::EMAIL, self::PASSWORD)
					),
					],
					content: json_encode($data),
					);
					$content = json_decode($response->getContent(), true);
					$this->assertEquals(201, $response->getStatusCode());
					$vacancyData = $this->em->getRepository(Vacancy::class)->findOneBy(['id' => $content['id']]);
					$this->assertEquals($data['name'], $vacancyData->getName());
					$this->assertEquals($data['workAddress'], $vacancyData->getWorkAddress());
				}
				
				public function testApiNewVacancyWrong()
				{
					$data = [
					'name' => '',
					'workAddress' => '',
					'description' => '',
					'specializations' => [
					1
					],
					];
					$response = $this->apiRequest(
					$this->router->generate('api_registration'),
					'POST',
					content: json_encode($data),
					);
					
					$this->assertEquals(400, $response->getStatusCode());
					$content = json_decode($response->getContent(), true);
					$this->assertEquals('Validation Failed', $content['form']['message']);
				}
			\end{lstlisting}  
			\caption{Модульное тестирование метода создания пользователя
				класса VacancyApiController}
			\label{unitvaca:image}
		\end{figure}
	\end{itemize} 
\end{enumerate}


Тестирование класса ResumeApiController 

\begin{enumerate}
	\item Тест метода создания вакансии(‘newResume‘):
	\begin{itemize}
		\item проверить, что при передаче корректных данных резюме
		сохраняется в базе данных;
		\item проверить, что при отсутствии обязательных данных возвращается ошибка валидации.
		\begin{figure}[H]
			\begin{lstlisting}[language=PHP]
				public function testApiNewResume()
				{
					$this->createUser();
					$this->em->flush();
					
					$data = ['specialization' => 1, 'firstName' => 'test', 'lastName' => 'test', 'workSchedule' => 1,'employmentType' => [1], 'residenceCity' => 1,'citizenship' => [1] ];
					
					$response = $this->apiRequest(
					$this->router->generate('api_resume_new'),
					'POST',
					headers: [
					'HTTP_AUTHORIZATION' => sprintf('Bearer %s', $this->getAuthToken(self::EMAIL, self::PASSWORD)),],
					content: json_encode($data),
					
					);
					$content = json_decode($response->getContent(), true);
					$this->assertEquals(201, $response->getStatusCode());
					$resumeData = $this->em->getRepository(Resume::class)->findOneBy(['id' => $content['id']]);
					$this->assertEquals($data['firstName'], $resumeData->getFirstName());
					$this->assertEquals($data['lastName'], $resumeData->getLastName());
				}
				
				public function testApiNewResumeWrong()
				{
					$this->createUser();
					$this->em->flush();
					
					$data = ['specialization' => '', 'firstName' => 'test', 'lastName' => 'test'];
					
					$response = $this->apiRequest(
					$this->router->generate('api_resume_new'),
					'POST',
					headers: [
					'HTTP_AUTHORIZATION' => sprintf('Bearer %s', $this->getAuthToken(self::EMAIL, self::PASSWORD)), ],
					content: json_encode($data));
					
					$this->assertEquals(400, $response->getStatusCode());
					$content = json_decode($response->getContent(), true);
					$this->assertEquals('Validation Failed', $content['form']['message']);
				}
			\end{lstlisting}  
			\caption{Модульное тестирование метода создания пользователя
				класса ResumeApiController}
			\label{resu:image}
		\end{figure}
	\end{itemize} 
\end{enumerate}

Модульное тестирование представляет собой критически важную составляющую процесса разработки программно-информационной системы. Данный подход способствует обнаружению ошибок в работе отдельных модулей системы, таких как функциональность добавления резюме, создания вакансий и процедуры регистрации пользователей. В ходе тестирования каждый метод и класс подвергается тщательной проверке на предмет корректности выполнения операций, включая обработку пограничных значений и некорректных входных данных. Такой подход существенно повышает качественные показатели и устойчивость разрабатываемого программного решения.

На рисунке~\ref{tests:image} представлен результат выполнения успешных тестов
для ключевых компонентов системы.

\begin{figure}[H]
	\centering
	\includegraphics[width=0.7\linewidth]{tests}
	\caption{Результат успешного тестирования}
	\label{tests:image}
\end{figure}

\subsection{Системное тестирование платформы}

Целью данного тестирования является проверка функциональности, надежности и производительности веб-платформы для поиска работы.
Системное тестирование позволяет выявить и устранить ошибки, а такжеоценить соответствие системы предъявленным требованиям.

Все тестовые сценарии должны проходить успешно без критических сбоев. Платформа должна обеспечивать стабильную и удобную работу пользователей в соответствии с установленными бизнес-процессами. Результаты тестирования будут использованы для повышения качества, надежности и пользовательского опыта сервиса по поиску работы.

\textbf{1) Запуск системы}

Описание: система должна запускаться без ошибок и отображать глав-
ное окно интерфейса.

На рисунке \ref{landing:image} представлена главная страница сайта.
\begin{figure}[H] 
\center{\includegraphics[width=1\linewidth]{landing}}
\caption{Главная страница}
\label{landing:image}
\end{figure}

\textbf{2) Авторизация пользователя}

Описание: Пользователь должен иметь возможность войти в систему с
использованием корректных учетных данных.

Ожидаемый результат: При вводе верного имени пользователя и пароля
осуществляется успешный вход, открывается главная рабочая область. При
вводе некорректных данных система отображает сообщение об ошибке.

На рисунке \ref{auth:image} представлена страница авторизации.

\begin{figure}[H]
\center{\includegraphics[width=1\linewidth]{auth}}
\caption{Страница авторизации}
\label{auth:image}
\end{figure}

\textbf{3) Проверка отображения данных авторизованного пользователя}

\begin{enumerate}
	\item Перейти на страницу \enquote{Профиль} и убедиться, что страница
	отображается корректно.
\end{enumerate}

На рисунке \ref{profile:image} представлена страница \enquote{Профиль}.

\begin{figure}[H]
	\center{\includegraphics[width=1\linewidth]{profile}}
	\caption{страница \enquote{Профиль} }
	\label{profile:image}
\end{figure}

\textbf{4) Добавление вакансии}
\begin{enumerate}
	\item Нажать кнопку \enquote{Создать вакансию}.
	\item Заполнить форму.
\end{enumerate}

На рисунке \ref{vacancy:image} представлена страница \enquote{Создание вакансии}.

\begin{figure}[H]
	\center{\includegraphics[width=1\linewidth]{vacancy}}
	\center{\includegraphics[width=1\linewidth]{vacancy2}}
	\caption{страница \enquote{Создание вакансии} }
	\label{vacancy:image}
\end{figure}

\textbf{5) Проверка отображения вакансий пользователя}
\begin{enumerate}
	\item Перейти на страницу \enquote{Мои вакансии}  и убедиться, что страница
	отображается корректно.
\end{enumerate}

На рисунке \ref{vacpers:image} представлена страница \enquote{Мои вакансии}.

\begin{figure}[H]
	\center{\includegraphics[width=1\linewidth]{vacpers}}
	\caption{ Cтраница \enquote{Мои вакансии}}
	\label{vacpers:image}
\end{figure}

\textbf{6) Просмотр вакансии}
\begin{enumerate}
	\item Нажать кнопку \enquote{Просмотреть}.
\end{enumerate}

На рисунке \ref{vacview:image} представлена страница просмотра вакансии.

\begin{figure}[H]
	\center{\includegraphics[width=1\linewidth]{vacview}}
	\caption{ Cтраница просмотра вакансии}
	\label{vacview:image}
\end{figure}

\textbf{7) Редактирование вакансии}
\begin{enumerate}
	\item Нажать кнопку \enquote{Редактировать}.
	\item Изменить данные вакансии.
	\item Нажать кнопку \enquote{Сохранить изменения}.
\end{enumerate}

На рисунке \ref{vacvedit:image} представлена страница редактирования вакансии.

\begin{figure}[H]
	\center{\includegraphics[width=1\linewidth]{vacedit}}
	\caption{ Cтраница редактирования вакансии}
	\label{vacvedit:image}
\end{figure}

\textbf{7) Удаление вакансии}
\begin{enumerate}
	\item Нажать кнопку \enquote{Удалить}.
	\item Подтвердить удаление.
\end{enumerate}

На рисунке \ref{vacdel:image} представлено окно удаления вакансии.

\begin{figure}[H]
	\center{
\includegraphics[width=0.7\linewidth]{vacdel}}
	\caption{ Окно удаления вакансии}
	\label{vacdel:image}
\end{figure}

На рисунке \ref{vacdelres:image} представлен результат удаления вакансии.

\begin{figure}[H]
	\center{\includegraphics[width=1\linewidth]{vacdelres}}
	\caption{ Результат удаления вакансии}
	\label{vacdelres:image}
\end{figure}

\textbf{8) Добавление компании}
\begin{enumerate}
	\item Нажать кнопку \enquote{Создать компанию}.
	\item Заполнить форму.
\end{enumerate}

На рисунке \ref{compadd:image} представлена страница создания компании.

\begin{figure}[H]
	\center{\includegraphics[width=0.7\linewidth]{compadd}}
	\caption{Страница создания компании}
	\label{compadd:image}
\end{figure}

\textbf{8) Просмотр компаний пользователя}
\begin{enumerate}
	\item Перейти на страницу \enquote{Мои компании}  и убедиться, что страница
	отображается корректно.
\end{enumerate}

На рисунке \ref{compview:image} представлена страница просмотра компаний пользователя.

\begin{figure}[H]
	\center{\includegraphics[width=0.7\linewidth]{compview}}
	\caption{Страница просмотра компаний пользователя}
	\label{compview:image}
\end{figure}


\textbf{9) Редактирование компании}
\begin{enumerate}
	\item Нажать кнопку \enquote{Редактировать}.
	\item Изменить данные.
\end{enumerate}

На рисунке \ref{compedit:image} представлена страница редактирования компании.

\begin{figure}[H]
	\center{\includegraphics[width=1\linewidth]{compedit}}
	\caption{Страница редактирования компании}
	\label{compedit:image}
\end{figure}


\textbf{10) Удаление компании}
\begin{enumerate}
	\item Нажать кнопку \enquote{Удалить}.
	\item Подтвердить удаление.
\end{enumerate}

На рисунке \ref{compdel:image} представлено окно удаления компании.

\begin{figure}[H]
	\center{
\includegraphics[width=0.7\linewidth]{compdel}}
	\caption{Окно удаления компании}
	\label{compdel:image}
\end{figure}

На рисунке \ref{compdelres:image} представлен результат удаления компании.

\begin{figure}[H]
	\center{\includegraphics[width=1\linewidth]{compdelres}}
	\caption{ Результат удаления компании}
	\label{compdelres:image}
\end{figure}

\textbf{11) Проверка отображения доступных резюме}
\begin{enumerate}
	\item Перейти на страницу поиска резюме и убедиться, что страница отоб-
	ражается корректно.
\end{enumerate}

На рисунке \ref{resumelist:image} представлена страница поиска резюме.

\begin{figure}[H]
	\center{\includegraphics[width=1\linewidth]{resumelist}}
	\caption{Страница поиска резюме}
	\label{resumelist:image}
\end{figure}

\textbf{12) Фильтрация}
\begin{enumerate}
	\item Перейти на страницу поиска резюме или вакансий.
	\item Выбрать фильтр.
	\item Нажать кнопку \enquote{Применить фильтры}.
\end{enumerate}

На рисунке \ref{filters:image} представлено окно фильтрации.

\begin{figure}[H]
	\center{\includegraphics[width=0.3\linewidth]{filters}}
	\caption{Окно фильтрации}
	\label{filters:image}
\end{figure}

На рисунке \ref{filtersres:image} представлен результат фильтрации.

\begin{figure}[H]
	\center{\includegraphics[width=1\linewidth]{filtersres}}
	\caption{ Результат фильтрации}
	\label{filtersres:image}
\end{figure}


\textbf{13) Добавление резюме}
\begin{enumerate}
	\item Нажать кнопку \enquote{Создать резюме}.
	\item Заполнить форму.
\end{enumerate}

На рисунке \ref{resumeform:image} представлена страница \enquote{Создание резюме}.

\begin{figure}[H]
	\center{\includegraphics[width=0.7\linewidth]{resumeform}}
	\center{\includegraphics[width=0.7\linewidth]{resumeform2}}
	\caption{Страница \enquote{Создание резюме} }
	\label{resumeform:image}
\end{figure}

\textbf{14) Проверка отображения резюме пользователя}
\begin{enumerate}
	\item Перейти на страницу \enquote{Мои резюме}  и убедиться, что страница
	отображается корректно.
\end{enumerate}

На рисунке \ref{resumepers:image} представлена страница \enquote{Мои резюме}.

\begin{figure}[H]
	\center{\includegraphics[width=1\linewidth]{resumepers}}
	\caption{ Cтраница \enquote{Мои резюме}}
	\label{resumepers:image}
\end{figure}

\textbf{15) Редактирование резюме}
\begin{enumerate}
	\item Нажать кнопку \enquote{Редактировать}.
	\item Изменить данные.
	\item Нажать кнопку \enquote{Сохранить изменения}.
\end{enumerate}

На рисунке \ref{resumeedit:image} представлена страница редактирования резюме.

\begin{figure}[H]
	\center{\includegraphics[width=0.6\linewidth]{resumeedit}}
	\caption{ Cтраница редактирования резюме}
	\label{resumeedit:image}
\end{figure}

\textbf{16) Удаление резюме}
\begin{enumerate}
	\item Нажать кнопку \enquote{Удалить}.
	\item Подтвердить действие.
\end{enumerate}

На рисунке \ref{resumedel:image} представлено окно удаления резюме.

\begin{figure}[H]
	\center{
\includegraphics[width=0.8\linewidth]{resumedel}}
	\caption{ Окно удаления резюме}
	\label{resumedel:image}
\end{figure}

\textbf{17) Отклик на вакансию}
\begin{enumerate}
	\item Нажать кнопку \enquote{Откликнуться}.
	\item Выбрать резюме.
	\item Нажать кнопку \enquote{Отправить}.
\end{enumerate}

На рисунке \ref{otklik:image} представлено окно отклика.

\begin{figure}[H]
	\center{\includegraphics[width=0.8\linewidth]{otklik}}
	\caption{ Окно отклика}
	\label{otklik:image}
\end{figure}

На рисунке \ref{otklikres:image} представлен результат отклика.

\begin{figure}[H]
	\center{\includegraphics[width=0.8\linewidth]{otklikres}}
	\caption{ Результат отклика}
	\label{otklikres:image}
\end{figure}


